\documentclass[11pt]{article}

\usepackage[margin=1in]{geometry}
\usepackage{lmodern}
\usepackage{amsmath,amssymb}
\usepackage{booktabs}
\usepackage[numbers,sort&compress]{natbib}
\usepackage[hidelinks]{hyperref}

\title{Which KAN Edge?\\
  \large Spline vs.\ Fourier vs.\ One-Qubit Variational Activations}
\author{
  Akhil Anand\\
  Arizona State University\\
  \texttt{aanan102@asu.edu}
}
\date{}

\begin{document}
\maketitle

\begin{abstract}
Kolmogorov--Arnold Networks place a learnable univariate function on every
edge~\citep{liu2024kan}. Prior fair comparisons find this design most
competitive on scientific formula-fitting, not as a general MLP
replacement~\citep{yu2024fairer}. We ask a narrower question: \emph{within
that regime, which edge basis is best?} We compare B-spline KANs, Fourier
KANs~\citep{xu2024fourierkan,mehrabian2024fourierkan}, and QKANs whose edge
nonlinearity is a one-qubit data-reuploading circuit with $\langle Z\rangle$
readout~\citep{jiang2025qkan,perezsalinas2020reuploading}. All three share the
same SiLU residual; an MLP is a sideline
reference~\citep{rumelhart1986backprop}. Models are trained in
JAX~\citep{eigenflow2026} on AI Feynman equations and special
functions~\citep{udrescu2019aifeynman} at label noise $0$ and $0.1$. We
report test RMSE against trainable parameter count, analytic forward FLOPs,
and wall-clock time. QKAN edges are exact $2\times 2$ statevector simulations,
not hardware; we do not claim quantum advantage.
\end{abstract}

\section{Introduction}
\label{sec:intro}

\section{Experimental Setup}
\label{sec:setup}

\section{Results}
\label{sec:results}

\section{Conclusion}
\label{sec:conclusion}

\bibliographystyle{plainnat}
\bibliography{references}

\end{document}
